\begin{thebibliography}{99}
\bibitem{rlc}
Cours electronique, \textit{Pont RLC},
\\
                \texttt{https://courselectronique.wordpress.com/etude-des-circuits/
                circuits-en-regime-transitoire/circuit-rlc-serie/}
\\
        Consulte le : 7 Janvier 2026
\bibitem{protocole}
        \\
                Protocole de laboratoire, \textit{Hepia, cours sur les ondes}
                \texttt{https://cyberlearn.hes-so.ch/pluginfile.php/697672/
                                mod\_resource/content/6/MT2\_S2\_main.pdf}

        \\
    \bibitem{mistral}
            Online Free LLM and Local LLM\textit{Mistral3 (online), Mistral:7b (local)},
        \textbf{Utilisé pour :}  
        \begin{itemize}
            \item \textbf{Introduction}
                : reformulation et orthographe
            \item \textbf{Théorie}
                : Mise en forme des équations
            \item \textbf{Description manipulations}
                : fautes d'orthographes
            \item \textbf{Résultats et analyse}
                : Mise en forme des tableaux et équations + vérification numérique $\Delta \gamma$
            \item \textbf{Conclusion}
                : Reformulation de syntaxe douteuse et fautes d'orthographes
\end{thebibliography}

\newpage

\section{Annexes}


\begin{figure}[H]
        \centering
            \includegraphics[scale=0.55]{ohwaw.png}
                \caption{Illustration d'ondes constructives et destructives.}
                    \label{fig:ohwaw}
\end{figure}
\newpage

\subsection{Ondes accoustiques}

\subsubsection{Data et graphiques}

\begin{center}
    \begin{tabular}{|c|c|c|c|c|}
        \hline
        \textbf{f} $\pm$ 1 $\backslash$ Hz &
        \textbf{$m 1$} $\pm$ 1 $\backslash$ mm &
        \textbf{$m_2$} $\pm$ 1 $\backslash$ mm &
        \textbf{$\Delta$m} $\pm$ 2 $\backslash$ mm &
        $\lambda$ $\pm$ 0.004 $\backslash$ m \\
        \hline
        900  & 99  & 291 & 192 & 0,384 \\
        1300 & 90  & 225 & 135 & 0,270 \\
        1700 & 63  & 166 & 103 & 0,206 \\
        2150 & 61  & 142 & 81  & 0,162 \\
        2550 & 42  & 110 & 68  & 0,136 \\
        3000 & 34  & 91  & 57  & 0,114 \\
        3500 & 32  & 81  & 49  & 0,098 \\
        4000 & 28  & 71  & 43  & 0,086 \\
        4500 & 30  & 68  & 38  & 0,076 \\
        5000 & 26  & 60  & 34  & 0,068 \\
        \hline
    \end{tabular}
    \captionof{table}{Mesures des interférences acoustiques (m $\rightarrow$ minimas)}
    \label{tab:konig_minimas}
\end{center}
    \newpage

\subsubsection{Détails calculs et propagation d'incertitudes}
\\Grâce à Eq. \ref{eq:celerite} et sachant que \(\lambda\) est obtenue par la differences entre positions des minimas à \(\pm \SI{1}{\milli\meter}\), l'incertitude sur \(\lambda\) est alors de \(\pm \SI{0.004}{\meter}\). L'incertitude absolue sur \( c \) est déjà obtenue par le fit du code python suivant :

\begin{figure}[h!]
    \centering
    \includegraphics[scale=0.6]{code_konig.png}
    \label{fig:code_konig}
    \caption{Script du fit de vitesse du son dans l'air}
\end{figure}


La longueur d'onde $\lambda$ est obtenue par double de m. Soit $\Delta \lambda$ = $2 \cdot \frac{\Delta m}{m}$.  
\newpage
\subsection{Ondes de surface circulaires}


\subsubsection{Data et graphiques}

\begin{center}
    \begin{tabular}{|c|c|c|c|c|c|}
        \hline
            \(f \pm 0.5\) Hz & \(D \pm 1\) $\backslash$ mm & \(N\) & \(\lambda \pm \Delta\lambda\) $\backslash$ m & \(\lambda_c \pm \Delta\lambda_c\) $\backslash$ m & \(c \pm \Delta c\) $\backslash$ m/s \\
        \hline
        14.8 & 40 & 2 & \(0.02000 \pm 0.00102\) & \(0.0120 \pm 0.0007\) & \(0.177 \pm 0.013\) \\
        18.0 & 53 & 3 & \(0.01767 \pm 0.00061\) & \(0.0106 \pm 0.0005\) & \(0.191 \pm 0.012\) \\
        20.9 & 48 & 3 & \(0.01600 \pm 0.00055\) & \(0.0096 \pm 0.0004\) & \(0.200 \pm 0.011\) \\
        31.0 & 47 & 4 & \(0.01175 \pm 0.00031\) & \(0.0070 \pm 0.0003\) & \(0.218 \pm 0.011\) \\
        40.6 & 60 & 5 & \(0.01200 \pm 0.00025\) & \(0.0060 \pm 0.0002\) & \(0.242 \pm 0.009\) \\
        49.6 & 60 & 6 & \(0.01000 \pm 0.00018\) & \(0.0051 \pm 0.0002\) & \(0.254 \pm 0.008\) \\
        61.1 & 45 & 5 & \(0.00900 \pm 0.00020\) & \(0.0045 \pm 0.0002\) & \(0.275 \pm 0.007\) \\
        70.1 & 49 & 6 & \(0.00817 \pm 0.00015\) & \(0.0042 \pm 0.0001\) & \(0.294 \pm 0.007\) \\
        79.2 & 50 & 7 & \(0.00714 \pm 0.00011\) & \(0.0037 \pm 0.0001\) & \(0.296 \pm 0.007\) \\
        \hline
    \end{tabular}
    \captionof{table}{Mesures des distances D, des fréquences f et des longueurs d'onde $\lambda = \frac{D}{N}$ avec N le nombre de mesures et $lambda_c$ la longueur d'onde corrigée en fonction du facteur de grossissement.}
    \label{tab:minimas_gamma}
\end{center}
\begin{figure}[h!]
    \centering
    \includegraphics[scale=0.5]{incertitudes_lambda_c.png}
    \label{fig:incert}
    \caption{Incertitude de $\lambda_c$ en fonction de la fréquence}
\end{figure}


\newpage


\begin{figure}[h!]
    \centering
    \includegraphics[scale=0.5]{IMG_4763.jpeg}
        \caption{Image pre-traitée, Distance entre excitateurs D = (0.113 +/- 0.001) m à fréquence f = (25.2 +/- 0.1) Hz (TO DEFINE) }
    \label{fig:4763_origin}
\end{figure}
\newpage
\begin{figure}[h!]
    \centering
    \includegraphics[scale=0.5]{4763_interferences_with_hyperboles.png}
    \caption{Image post-traitement avec superposition}
    \label{fig:4763_pixels}
\end{figure}

\begin{figure}[h!]
    \centering
    \includegraphics[scale=0.5]{4763_interferences_with_scaled_hyperboles.png}
    \caption{image post-traitement sans superposition a l'echelle}
    \label{fig:4763_scaled}
\end{figure}

\newpage

\begin{figure}[h!]
    \centering
        \includegraphics[scale=0.5]{IMG_4771.jpeg}
    \caption{Image pre-traitée, Distance entre excitateurs D=0.041m à fréquence f = (19.6 +/- 0.1) Hz}
    \label{fig:4771_origin}
\end{figure}
\newpage
\begin{figure}[h!]
    \centering
        \includegraphics[scale=0.5]{4771_interferences_with_hyperboles.png}
        \caption{Image post-traitement avec superposition. Distance entre excitateurs D=0.041m à fréquence f = (19.6 +/- 0.1) Hz}
        \label{fig:4771}
\end{figure}
\begin{figure}[h!]
    \centering
        \includegraphics[scale=0.5]{4771_interferences_with_scaled_hyperboles.png}
        \caption{Image post-traitement sans superposition a l'echelle.  Distance entre excitateurs D=0.041m à fréquence f = (19.6 +/- 0.1) Hz}
        \label{fig:4771_scaled}
\end{figure}
\newpage


\begin{figure}[h!]
    \centering
    \includegraphics[scale=0.5]{IMG_4774.jpeg}
    \caption{Image pre-traitée, Distance entre excitateurs D=0.041 +/- 0.001 m à fréquence d'excitation f = (25.2 +/- 0.1) Hz}
    \label{fig:4774_origin}
\end{figure}
\newpage
\begin{figure}[h!]
    \centering
        \includegraphics[scale=0.5]{4774_interferences_with_hyperboles.png}
        \caption{Image post-traitement avec superposition Distance entre excitateurs D=0.041 +/- 0.001 m à fréquence d'excitation f = (25.2 +/- 0.1) Hz}
        \label{fig:4774}
\end{figure}
\begin{figure}[h!]
    \centering
        \includegraphics[scale=0.5]{4774_interferences_with_scaled_hyperboles.png}
        \caption{Image post-traitement sans superposition a l'echelle. Distance entre excitateurs D=0.041 +/- 0.001 m à fréquence d'excitation f = (25.2 +/- 0.1) Hz}
        \label{fig:4774_scaled}
\end{figure}
\newpage


\subsubsection{Details de calculs et propagationd d'incertitudes}

\textbf{Incertitudes mesure tension superficielle}

Les longueurs d'onde sans le grossissement sont mesurées par rapport au quotient du nombre mesurées de ces dernières (N) et la distance mesurée D. Soit :
\[
    \Delta \lambda = \frac{\Delta D}{D \cdot N}
\]
avec $\Delta D = 0.001 m$.


La véritable longueur d'onde ($\Lambda_c$) est alors est mesurée avec un facteur de grossissement dont l'incertitude est calculée à partir des incertitudes sur \(x_1\) et \(x_2\). $\Delta \lambda_c$ depend alors du rapport entre le grossissement et de $\lambda$. Soit :

\[
    \Delta \lambda_{\text{c}} = \lambda_{\text{c}} \cdot \sqrt{\left(\frac{\Delta x_1}{x_1}\right)^2 + \left(\frac{\Delta x_2}{x_2}\right)^2 + \left(\frac{\Delta \lambda}{\lambda}\right)^2
\]

où \(\Delta x_1 = \Delta x_2 = \SI{0.001}{\meter}\) pour un grossissement g(x, y) avec x $\approx y$ où x1 = 55 $\pm 1mm$ et sa projection g(x1) = 91 $\pm 1$ mm, y1 = 49 $\pm 1 mm$ et g(y1) = 82 $\pm 1 mm$.

\\

La période \(T\) est obtenue à partir de la fréquence \(f\), avec une incertitude \(\Delta f = \SI{0.5}{\hertz}\) :
\[
    \Delta T = \frac{\Delta f}{f^2}
\]
\\
La vitesse de phase \(c\) est donnée par :
\[
    c = \frac{\lambda_{\text{corr}}}{T}
\]
\\
L'incertitude sur \(c\) est calculée par propagation des incertitudes :

\[
    \Delta c = c \cdot \sqrt{\left(\frac{\Delta \lambda_{\text{c}}}{\lambda_{\text{c}}\right)^2 + \left(\frac{\Delta T}{T}\right)^2}
\]

\begin{equation}
    \Delta \gamma = \sqrt{
    \left( \frac{\partial \gamma}{\partial c} \Delta c \right)^2 +
    \left( \frac{\partial \gamma}{\partial \lambda} \Delta \lambda \right)^2 +
    \left( \frac{\partial \gamma}{\partial d} \Delta d \right)^2
    }
    \label{eq:incertitude_gamma}
\end{equation}

où les dérivées partielles sont calculées comme suit :

\begin{align}
    \frac{\partial \gamma}{\partial c} &= \frac{\rho}{\pi} \cdot \frac{c}{\tanh\left(\frac{2\pi d}{\lambda}\right)} \cdot \lambda, \\
    \frac{\partial \gamma}{\partial \lambda} &= \frac{\rho}{2\pi} \left( \frac{c^2}{\tanh\left(\frac{2\pi d}{\lambda}\right)} - \frac{g \lambda}{2\pi} \right) + \frac{\rho}{2\pi} \left( \frac{4\pi d c^2}{\lambda^2 \sinh^2\left(\frac{2\pi d}{\lambda}\right)} + \frac{g \lambda}{2\pi} \right), \\
    \frac{\partial \gamma}{\partial d} &= -\frac{\rho c^2}{\pi \sinh\left(\frac{4\pi d}{\lambda}\right)} \cdot \frac{4\pi}{\lambda}.
\end{align}

\textbf{Code python du fit de gamma}

\begin{figure}[h!]
    \centering
    \includegraphics[scale=0.6]{code_gamma.png}
    \label{fig:code_gamma}
    \caption{Script du fit de la tension superficielle $\gamma$}
\end{figure}

\subsubsection{Interferences : codes python des hyperboles}
\begin{figure}[h!]
    \centering
    \includegraphics[scale=0.6]{code_hyper_scaled.png}
    \label{fig:code_hyper}
    \caption{Script de la superposition des hyperboles et des maximas à l'echelle}
\end{figure}

\begin{figure}[h!]
    \centering
    \includegraphics[scale=0.8]{code_hyper.png}
    \label{fig:code_hyper}
    \caption{Partie du script afin d'afficher l'image post-traitement superposée aux hzperboles}
\end{figure}




\subsubse
